\section{Summary of Contributions}
\label{sec:conclusion-summary}

\begin{itemize}
    \item \textbf{Single-vector ERE backend} (Chapter~\ref{chap:ere}):
    consolidates two-layer metadata into one bit-vector, cutting metadata
    size by $24\%$.
    \item \textbf{Adaptive bucket-search} (Chapter~\ref{chap:ere}):
    drops $59$ bits of per-key metadata and speeds queries by
    $13$--$30\times$ at the bucket sizes that arise at per-SSTable scale.
    \item \textbf{Truncation hash} (Chapter~\ref{chap:hash}):
    saves $\log_2 \mathcal{L}$ bits per key on sparse inputs,
    going below the Goswami bound on distributions where keys are spread
    across the universe.
    \item \textbf{Seg-ARE} (Chapter~\ref{chap:hybrid}):
    a single-pass segmentation that routes dense clusters to exact-mode
    ERE sub-filters and sparse tails to the truncation hash, going below
    the bound on structured distributions.
\end{itemize}
The first two improvements bring the Goswami-line filter closer to the
lower bound on adversarial inputs; the truncation hash and Seg-ARE each
cross below it on inputs where the respective structure is present.

In absolute terms, the combined filter achieves the following on
standard SOSD benchmarks at $\mathrm{FPR} \leq 10^{-3}$:

\begin{itemize}
    \item \textbf{3.7 bits per key} on Wikipedia timestamps --- less
    than a third of any competitor at the same error rate.
    \item \textbf{7.7--9.2 bits per key} on clustered synthetic keys,
    versus $12$--$19$ for the next-best filter.
    \item \textbf{150--420\,ns} per query across all four SOSD
    real-world datasets and query lengths
    $\mathcal{L} \in \{1,\,16,\,1024\}$
    (Table~\ref{tab:latency}).
\end{itemize}

\section{When Seg-ARE Wins, Ties, and Loses}
\label{sec:conclusion-where}

The empirical picture across the four SOSD distributions falls out of
the per-cluster vs universe-wide Elias--Fano decomposition of
Section~\ref{sec:hybrid-ef}:

\begin{itemize}
    \item \textbf{Wiki --- clearest win.} Wikipedia edit timestamps
    are heavily concentrated on the right side of the universe
    (Figure~\ref{fig:sosd-histograms}). The segmentation isolates the dense parts
    cleanly; per-cluster EF encodes each tight region in far fewer
    bits than a universe-wide EF could.
    Seg-ARE reaches $\mathrm{FPR} = 0$ at $3.7$~BPK against Grafite's
    $4.6$~BPK at the lossless threshold
    (Figure~\ref{fig:seg-wiki-L128}, Table~\ref{tab:bpk}, $\mathcal{L}=16$ block).
    \item \textbf{Facebook --- tied.} Facebook user IDs occupy a
    relatively small $37$-bit universe with no large gaps. There is
    little for segmentation to exploit, and universe-wide EF is
    already compact.
    Seg-ARE and Grafite both reach FPR $= 0$ at $\sim\!11$ BPK;
    Seg-ARE matches but does not improve over the Grafite-with-fallback baseline.
    \item \textbf{OSM --- no advantage.} OpenStreetMap cell IDs are
    spread thinly across the full $64$-bit universe, with median gaps
    of order $10^7$. The DBSCAN-style detector finds essentially no
    clusters above the $\delta = 256\mathcal{L}/\varepsilon$ window,
    so most keys go to the truncation-hash fallback and Seg-ARE
    degrades to truncation. Honest hashing is the only option on this
    geometry; no segmentation primitive helps.
\end{itemize}

Seg-ARE is therefore not a universal improvement over Grafite.
It is a distribution-aware option that activates when the data has the
cluster-and-gap geometry common in many real LSM key sets
(Wiki-like clustered timestamps) and matches Grafite on the rest.

\section{Limitations}
\label{sec:conclusion-limits}

The Seg-ARE design carries three honest limitations.

\textbf{Detector sensitivity.} The forward sweep depends on
$\delta = \texttt{minPts} \cdot \mathcal{L} / (2\varepsilon)$ and on
the fixed choice $\texttt{minPts} = 512$
(Chapter~\ref{chap:hybrid}). On clusters smaller than $\texttt{minPts}$
keys the detector misses real density and the win is lost.
We have not investigated adaptive $\texttt{minPts}$ tuning.

\textbf{Per-segment overhead.} Each segment carries $\sim\!700$~bits
of fixed metadata (cluster bounds, sub-filter header, pointer). At
$\texttt{minPts} = 512$ this is $\sim\!1.4$~BPK of fixed overhead. On
inputs that produce many small clusters the overhead can outweigh the
per-cluster EF savings.

\textbf{Static setting.} The construction takes a fixed key set;
segments are not maintained under inserts or deletes. Dynamic key
sets, where the cluster structure itself evolves, require the
keepsake-box scheme of Memento~\citep{eslami2024memento} or an
equivalent online segmentation mechanism.

\section{Future Work}
\label{sec:conclusion-future}

Three directions stand out.

\textbf{Adaptive segmentation parameters.} The window and density
thresholds are currently fixed at construction time from a single
\texttt{minPts} value.
Choosing them per LSM level (top levels have
fewer keys but tighter clusters; bottom levels are denser but flatter)
could close the gap on OSM-like distributions where the global
$\delta$ is too coarse.

\textbf{Combining segmentation with learned models.}
The CDF-fitting approach of SNARF~\citep{vaidya2022snarf} solves the universe-mapping
problem that hurts Grafite on small-universe inputs, but loses the
worst-case guarantee. Running a segmenter first, then a CDF model on
each cluster, would localise the model's assumptions and might recover
the guarantee on the clusters individually.

\textbf{Dynamic Seg-ARE.} Combining the segmentation primitive with
Memento's dynamic keepsake-box machinery would extend the
distribution-aware win condition from static SSTables to mutable
indexes (B-trees, single-global-filter LSM-trees, mixed
transactional/analytical workloads).
